\section{A wider fixed-count band with the original central rate}
\label{sec:rate-preserving-band}
The multiscale stopping estimate can be combined with damped cubic
unsmoothing without altering the actual return mark. We choose the
outer radius so that the integrated stopping error has the same power
as the original marked central estimate. Thus the wider band no longer
requires a slower probability-decay budget for a single marked event.

Set
\begin{equation}\label{eq:rate-preserving-scales}
 \varepsilon_\natural=\frac{67}{1400},\qquad
 V_n^\natural=n^{67/1400},\qquad
 B_n^\natural=n^{-633/1400}.
\end{equation}
The smooth Fourier support radii are $2V_n^\natural$ after rescaling
and $2B_n^\natural$ in physical frequency. The preceding radius
$n^{1/21}$ is strictly smaller than $V_n^\natural$. All statements
below are at the prescribed return count $n$, not averaged in $n$.

\begin{theorem}[Rate-preserving prescribed-count annulus]
\label{thm:rate-preserving-annulus}
Uniformly in $R\in I$, $0\le k\le n$ and the marked bounded-BV
insertion $a$, for all sufficiently large $n$,
\begin{align}
 &n^2\int_{2n^{-99/200}\le|z|\le2B_n^\natural}
                     |\Phi^{[k],a}_{n,R}(z)|\dd z
       \le C\left(M_an^{-3/280}+V_an^{-983/350}\right),
       \label{eq:rate-preserving-annulus}\\
 &\int_{|v|\le2V_n^\natural}
       |C^{[k],a}_{n,R}(v)-\alpha_a e^{-v^{\mathsf T}D_Rv/2}|\dd v
       \le C\left(M_an^{-3/280}\sqrt{\log(2+n)}
                           +V_an^{-9/175}\right).
       \label{eq:rate-preserving-central}
\end{align}
The original probability corresponds to $a=s_R$. The transform with a
general insertion is still integrated against $\nu|_{Y_R^*}$, with
mass $\alpha_a=\int a\dd\nu$.
\end{theorem}
\begin{proof}
Recenter at the actual mark. Equation~\eqref{eq:quarter-characteristic-stopping}
compares the original transform with
$I^a_{m_k,m_{n-k},R}(v/\sqrt n)$ at cost $CM_a|v|n^{-1/4}$.
The deterministic length satisfies $m_k+m_{n-k}=n/c_*+O(1)$.
Apply Theorem~\ref{thm:damped-cubic-unsmoothing} with
\[
 Q=19,\qquad \delta=\tfrac14 n^{-1/5},\qquad
 \epsilon=\tfrac14 n^{-3},\qquad |v|\le2n^{67/1400}.
\]
The coarse scale determines the collision perturbation disk, while
only the insertions use the fine scale. For $z=v/\sqrt n$, the two
required quantities satisfy
\begin{align}
 |z|\delta^{-2}&=O(n^{-73/1400}),\notag\\
 |z|^{18}\delta^{-40}&=O(n^{-97/700}).
 \label{eq:rate-preserving-damping-domain}
\end{align}
Hence the damping theorem applies for all sufficiently large $n$.
On the annulus $2n^{1/200}\le|v|\le2V_n^\natural$, the damped
polynomial term is smaller than every fixed inverse power of $n$:
its prefactor is polynomial and its exponential is at most
$e^{-c n^{1/100}}$.

We display every nondamping integrated exponent. Four-dimensional
volume and the largest frequency factor give respectively
\begin{align}
 \epsilon V_a(V_n^\natural)^4
      &\le C V_an^{-983/350},\notag\\
 M_a\epsilon(1+\sqrt n V_n^\natural)^3(V_n^\natural)^4
      &\le C M_an^{-233/200},\notag\\
 M_a\delta^2L_\delta^2(V_n^\natural)^8
      &\le C M_an^{-3/175}[\log(2+n)]^2,\notag\\
 M_an^{-1}\delta L_\delta^3(V_n^\natural)^8
      &\le C M_an^{-143/175}[\log(2+n)]^3,\notag\\
 M_an^{-1/4}(V_n^\natural)^5
      &= M_an^{-3/280}.
 \label{eq:rate-preserving-exponents}
\end{align}
The quartic residual exponent exceeds $3/280$ by $9/1400$, so its
logarithms are absorbed by this strict power margin. Every other
supremum-norm error also decays faster than the last term. Since
$v=\sqrt n z$ has Jacobian $\dd v=n^2\dd z$ in four dimensions,
these estimates prove \eqref{eq:rate-preserving-annulus} with the
raw factor $n^2$ and no division by annular volume.

On $|v|\le2n^{1/200}$, use the sharper inherited
Theorem~\ref{thm:marked-return-band}. On the rest of the new ball,
use \eqref{eq:rate-preserving-annulus} and the uniformly elliptic
Gaussian tail. Since $983/350>9/175$, their sum gives
\eqref{eq:rate-preserving-central}. This does not replace the old
small-ball theorem by a weaker one.
\end{proof}

\begin{remark}[The choice of radius]
The equality
\[
 \frac14-5\varepsilon_\natural=\frac3{280}
\]
selects the radius for which this stopping estimate retains the original
central power. It is a balance of the displayed error bounds, not an
assertion of optimality or an obstruction to another method. With a
smaller positive rate, the preceding construction still reaches every
fixed $1/200<\varepsilon<1/20$. Neither statement covers the rescaled
endpoint $n^{1/10}$, corresponding to physical frequency $n^{-2/5}$.
\end{remark}

\begin{corollary}[The original rare-event budget on the wider band]
\label{cor:rate-preserving-same-event}
For the unchanged marked-state event $A_{n,k,R}$, probability $p_{n,R}$,
and indicator variation $V_{n,R}$ of
Corollary~\ref{cor:all-same-event-moments},
\begin{align}
 &\int_{|v|\le2V_n^\natural}
 \left|\mathbb E_{\nu_R^*}
       [e^{iv\cdot U_{n,R}/\sqrt n}\mid A_{n,k,R}]
                               -e^{-v^{\mathsf T}D_Rv/2}\right|\dd v\notag\\
 &\hspace{18mm}\le\frac C{p_{n,R}}
        \left(n^{-3/280}\sqrt{\log(2+n)}+V_{n,R}n^{-9/175}\right).
 \label{eq:rate-preserving-same-event}
\end{align}
Thus $p_{n,R}\ge cn^{-\beta}$ and $V_{n,R}\le Cn^\kappa$ suffice
when $\beta<3/280$ and $\beta+\kappa<9/175$. Under these conditions,
every fixed conditional polynomial moment also has the Gaussian limit.
\end{corollary}
\begin{proof}
Use $a=c_*^{-1}\one_{E_{n,R}}$ in
\eqref{eq:rate-preserving-central}. Its mass is the exact denominator
$p_{n,R}$. Division proves the bound without changing the event.
The stated inequalities also imply the moment conditions in
Corollary~\ref{cor:all-same-event-moments}.
\end{proof}

\subsection{The raw identity with the new kernel}
Define
\[
 \chi_n^\natural(\omega)=\chi(\omega/B_n^\natural),\qquad
 K_n^\natural=\mathcal F^{-1}\chi_n^\natural.
\]
For a marked insertion additionally satisfying the finite-record
admissibility in Definition~\ref{def:finite-record-weight}, use the
same weight $w_{n,k,R}=c_*a\circ(F_R^*)^k$ and the same exact finite
extraction $\mu^{a,L}=E^{a,L}+Q^{a,L}$ as in
Theorem~\ref{thm:sharp-weighted-raw-inversion}. Set, with the new
kernel throughout,
\begin{align}
 \mathcal E^{\natural,a,L}_{n,k,R,A}
 &=\mathop{\rm ess\,sup}_{\mathcal W_{n,R,A}}
                 |e^{a,L}-K_n^\natural*E^{a,L}|,\notag\\
 \mathcal C^{\natural,a,L}_{n,k,R}(B)
 &=\int_{\T^3\times[-B,B]}
       |1-\chi_n^\natural(\omega)|\,|\widehat Q^{a,L}(\omega)|\dd\omega.
 \label{eq:rate-preserving-raw-terms}
\end{align}
In particular the edge correction is not copied from a different
cutoff. It may still be infinite before its germs are estimated.

\begin{corollary}[Exact weighted raw budget with the retained central rate]
\label{cor:rate-preserving-raw-budget}
Let $L_n=\lceil\lambda n\rceil$ with $\lambda>c_*^{-1}+1$.
For $B\ge1$, $P>0$, and $y=((\ell,t)-n\bar G_R)/\sqrt n$,
\begin{align}
 &\mathop{\rm ess\,sup}_{\mathcal W_{n,R,A}}
        |n^2p_{n,R}^{[k],a}(\ell,t)-\alpha_a\mathfrak g_R(y)|\notag\\
 &\quad\le C\left(M_an^{-3/280}\sqrt{\log(2+n)}+V_an^{-9/175}\right)
      +CM_ae^{-c n^{67/700}}+C_{A,\lambda,P}M_an^{-P}\notag\\
 &\qquad+n^2\mathcal E^{\natural,a,L_n}_{n,k,R,A}
     +(2\pi)^{-4}n^2\mathcal C^{\natural,a,L_n}_{n,k,R}(B)
     +\frac{n^2}{\pi B}A_2(n,L_n,R,w_{n,k,R}).
 \label{eq:rate-preserving-raw-budget}
\end{align}
The inequality has the extended-real interpretation when needed.
Both cutoff radii are as in \eqref{eq:rate-preserving-scales}.
\end{corollary}
\begin{proof}
On the central count labels the high-count tail has no density
contribution. The same exact support and absolute-residual-inversion
argument gives
\begin{align*}
 p_{n,R}^{[k],a}={}&K_n^\natural*\mu_{n,R}^{[k],a}
       +(e^{a,L_n}-K_n^\natural*E^{a,L_n})\\
 &+\mathcal F^{-1}[(1-\chi_n^\natural)\widehat Q^{a,L_n}]
       -K_n^\natural*\mathsf T^{a,L_n}.
\end{align*}
Use \eqref{eq:rate-preserving-central} for the first term. The missing
Gaussian frequencies have $|v|\ge V_n^\natural$ and give the stated
exponential tail. Count separation is unchanged, but the kernel scale
is new. Its Schwartz estimate is
\[
 n^2\sup_{\mathcal W_{n,R,A}}|K_n^\natural*\mathsf T^{a,L_n}|
       \le C_{A,\lambda,J}M_a
                     n^{67/350-767J/1400}.
\]
Indeed $n^2(B_n^\natural)^4=n^{67/350}$ and
$nB_n^\natural=n^{767/1400}$. Choose $J$ for the prescribed power
$P$. Splitting the residual inverse at roof frequency $B$ gives
exactly the last two terms, by the inherited absolute second-derivative
far-tail estimate. The local edge term stays in its combined form.
\end{proof}

\paragraph{What remains beyond this band.}
The new fixed-count estimate improves the physical transform, rather
than an average or a compressed operator. The outer region beginning
at $2n^{-633/1400}$, compact nonzero and peripheral return frequencies,
and growing roof frequencies still enter the complementary residual
integral. The long-time bound for $A_2$ at $L_n\asymp n$ and the local
edge correction remain distinct terms in
\eqref{eq:rate-preserving-raw-budget}. A relative comparison with an
exact physical-time event is not performed. The new all-moment and
same-event conclusions preserve the actual indicator and its exact
probability; they are not a substitute for those raw estimates.
